\documentclass[10pt]{article}
\usepackage[margin=0.62in]{geometry}
\usepackage{amsmath,amssymb,bm}
\usepackage{booktabs}
\usepackage{siunitx}
\usepackage{xcolor}
\usepackage{tcolorbox}
\usepackage{microtype}
\usepackage{enumitem}
\setlength{\parindent}{0pt}
\setlength{\parskip}{3pt}
\pagestyle{empty}
\sisetup{detect-all}
\definecolor{navy}{RGB}{28,62,82}
\definecolor{light}{RGB}{242,246,248}
\newcommand{\vect}[1]{\bm{#1}}

\begin{document}

{\Large\bfseries Measurement-updated Stokes prediction of $E_4$ visibility}\\[-1pt]
{\small Hybrid MZI: comparison of isolated-state polarization measurements with a swept-PAX power measurement}

\vspace{4pt}
\hrule
\vspace{5pt}

\textbf{1. Minimal update to the ideal model.}
The ideal path-$E_4$ field contains two coherent contributions whose relative phase is controlled by $\phi_1$. Replacing the ideal orthogonal basis states by the \emph{measured} normalized polarization states $\lvert A\rangle$ and $\lvert B\rangle$ gives
\begin{equation}
E_4(\phi_1)=\sqrt{P_A}\,e^{i\phi_1}\lvert A\rangle+\sqrt{P_B}\,\lvert B\rangle .
\end{equation}
The corresponding intensity is
\begin{equation}
I_4(\phi_1)=P_A+P_B+2\sqrt{P_AP_B}\,\Re\!\left[e^{i\phi_1}\langle B\vert A\rangle\right],
\end{equation}
so the theoretical fringe visibility is
\begin{equation}
\boxed{V_{\rm th}=\frac{2\sqrt{P_AP_B}}{P_A+P_B}\,\left|\langle A\vert B\rangle\right|}.
\end{equation}
For fully polarized states, the Jones-state overlap is obtained directly from the normalized Stokes vectors,
\begin{equation}
\left|\langle A\vert B\rangle\right|=
\sqrt{\frac{1+\vect{s}_A\!\cdot\!\vect{s}_B}{2}}.
\end{equation}

\textbf{2. Numerical substitution from the realigned isolated-state measurement.}
The measured states and powers were
\begin{align*}
\vect{s}_A&=(+0.990466,-0.051196,+0.127887), & P_A&=\SI{0.342329}{mW},\\
\vect{s}_B&=(-0.989706,-0.092665,-0.109066), & P_B&=\SI{0.417919}{mW}.
\end{align*}
Therefore
\begin{align*}
\vect{s}_A\cdot\vect{s}_B&=-0.989475,\\
\left|\langle A\vert B\rangle\right|&=\sqrt{\frac{1-0.989475}{2}}=0.072544,\\
\frac{2\sqrt{P_AP_B}}{P_A+P_B}&=0.995045,
\end{align*}
which gives
\begin{equation}
\boxed{V_{\rm th}=(0.995045)(0.072544)=0.07219\approx\mathbf{7.22\%}.}
\end{equation}
Thus the measured power imbalance is a small correction; the dominant predicted contrast comes from the measured nonorthogonality of the two polarization states.

\textbf{3. Visibility measured directly by the PAX during the $\phi_1$ sweep.}
The swept acquisition recorded $S_0=P_{\rm PAX}$ simultaneously with the Stokes trajectory. After excluding the initial $\sim\SI{1}{s}$ settling region, the observed power span was approximately
\[
P_{\min}=\SI{0.7264}{mW},\qquad P_{\max}=\SI{0.8615}{mW},
\]
corresponding to a literal-extrema visibility of $V_{\rm raw}=8.51\%$. Because the PAX samples only about 12 points per actuator period, cycle-wise extrema tend to miss the true extrema: the repeated-cycle mean is $6.60\%\pm0.61\%$. A pooled 2--98\% power envelope gives $V_{\rm PAX}=7.23\%$.

\begin{center}
\begin{tabular}{lcc}
\toprule
Quantity & Visibility & Interpretation \\
\midrule
Stokes/Jones prediction & \textbf{7.22\%} & No fit; isolated measured states and powers \\
PAX pooled 2--98\% envelope & \textbf{7.23\%} & Robust repeated-sweep estimate \\
PAX cycle-wise mean & 6.60\% $\pm$ 0.61\% & Sparse phase sampling lowers extrema \\
PAX literal global extrema & 8.51\% & Most sensitive to individual samples \\
\bottomrule
\end{tabular}
\end{center}

\begin{tcolorbox}[colback=light,colframe=navy,boxrule=0.6pt,arc=1.5pt,left=5pt,right=5pt,top=4pt,bottom=4pt]
\textbf{Result.} The simplest measurement-updated Jones model predicts $V_{\rm th}=7.22\%$. The independently swept PAX power data place the observed contrast in the same $\sim7\%$ range; the robust pooled estimate is $7.23\%$. Within the sampling limitations of the PAX sweep, the measured polarization nonorthogonality is therefore quantitatively sufficient to account for the visibility observed by the PAX in $E_4$.
\end{tcolorbox}

\end{document}
